\chapter{Requirement Analysis}

\section{Functional Requirements}

Functional requirements define the core operations and capabilities of the \textbf{SafeSurf} system. The system accepts input URLs, processes them through a multi-layered machine learning pipeline, and produces phishing detection results. The goal is to ensure accurate, real-time classification with minimal latency while maintaining a simple and user-friendly interface.

\subsection{Core Functional Behavior of SafeSurf}

\subsubsection*{User Interface \& Interaction}

\begin{itemize}
    \item SafeSurf provides a React.js-based web application that allows users to manually enter and analyze URLs, along with a Chrome extension that continuously monitors URLs in the background during browsing sessions.
    
    \item Users are able to input URLs directly through the web dashboard or rely on the Chrome extension to automatically scan URLs without requiring manual intervention.
    
    \item The system returns a detailed response that includes the classification result (\textbf{Safe, Phishing, or Invalid Domain}) along with a clearly defined risk level categorized as \textbf{Low, Medium, or High}.
    
    \item The user interface is designed to be simple, intuitive, and responsive, ensuring that users can easily understand results without being exposed to internal model complexities or raw probability values.
\end{itemize}

\subsubsection*{URL Processing Pipeline}

\begin{itemize}
    \item The system accepts a URL as input and performs normalization steps such as removing unnecessary ports, handling trailing slashes, and ensuring consistent formatting across different inputs.
    
    \item Input validation and sanitization mechanisms are applied to detect malformed or potentially harmful inputs before further processing.
    
    \item A DNS Guard layer performs domain resolution checks, identifies non-existent domains (NXDOMAIN), and filters out invalid URLs before they reach the machine learning pipeline.
    
    \item This preprocessing stage reduces unnecessary computation and improves system efficiency by ensuring that only valid URLs proceed to feature extraction and prediction.
\end{itemize}

\subsubsection*{Feature Extraction Engine}

\begin{itemize}
    \item The system extracts a total of 111 features from the input URL, covering lexical characteristics, network-level information, and textual patterns.
    
    \item Lexical features include properties such as URL length, structure, number of special characters, and token distribution within the URL.
    
    \item Network features include DNS records (A, NS, MX), WHOIS registration details, domain age, and SSL certificate-related information.
    
    \item NLP-based features are generated using a Bag-of-Words representation to capture textual patterns within URLs.
    
    \item Feature extraction is implemented using asynchronous and parallel processing techniques (such as \texttt{asyncio.gather}) to reduce latency and efficiently handle network-bound operations.
\end{itemize}

\subsubsection*{Machine Learning Detection Engine}

\begin{itemize}
    \item The system uses an ensemble of machine learning models including Logistic Regression, Linear Support Vector Machine (LinearSVC), Random Forest, Histogram Gradient Boosting, and a calibrated XGBoost classifier.
    
    \item A separate NLP-based Logistic Regression model is used to analyze textual characteristics of URLs independently of infrastructure-based features.
    
    \item The outputs from these models are combined using a weighted soft-voting approach, where infrastructure-based predictions contribute 65\% and text-based predictions contribute 35\% to the final result.
    
    \item This hybrid approach improves overall prediction accuracy by leveraging both structural and textual characteristics of URLs.
\end{itemize}

\subsubsection*{Real-Time Detection \& Response}

\begin{itemize}
    \item The system processes incoming requests in real time using a FastAPI backend that supports asynchronous request handling and efficient concurrency.
    
    \item The API returns a structured JSON response containing the classification result, risk level, and additional metadata derived from the prediction pipeline.
    
    \item Risk levels are determined using predefined probability thresholds, where values above a certain limit (e.g., 0.85) are categorized as high risk.
    
    \item The Chrome extension consumes this response to notify users and prevent access to high-risk websites during browsing.
\end{itemize}

\subsubsection*{System Maintainability}

\begin{itemize}
    \item The system follows a modular architecture in which machine learning models, feature extraction components, and API layers are separated for easier updates and maintenance.
    
    \item This modular design allows independent updates to models and feature pipelines without affecting the overall system.
    
    \item The system supports version control using Git and allows retraining and calibration of models as new data becomes available.
\end{itemize}

\subsubsection*{Software Requirements}

\begin{itemize}
    \item \textbf{Frontend:}
    \begin{itemize}
        \item React.js for the responsive web interface \cite{ref9}.
        \item Chrome Extension APIs for background script processing.
    \end{itemize}
    
    \item \textbf{Backend:}
    \begin{itemize}
        \item FastAPI (Python 3.9+) for core API architecture and machine learning pipeline integration \cite{ref10}.
    \end{itemize}
    
    \item \textbf{Machine Learning \& Data Processing:}
    \begin{itemize}
        \item Scikit-learn (VotingClassifier, NLP Models) \cite{ref11}.
        \item XGBoost (Infrastructure Intelligence) \cite{ref12, ref13}.
        \item NumPy and Pandas for data manipulation.
    \end{itemize}
    
    \item \textbf{Deployment:}
    \begin{itemize}
        \item Docker for multi-stage containerization \cite{ref14}.
    \end{itemize}
\end{itemize}

\subsubsection*{Hardware Requirements}

\begin{itemize}
    \item \textbf{Minimum Runtime Requirements:}
    \begin{itemize}
        \item Processor: Intel i5 / Ryzen 5 or equivalent.
        \item RAM: 8 GB (16 GB recommended).
    \end{itemize}
    \item \textbf{Training Environment Recommended:}
    \begin{itemize}
        \item RAM: 16 GB+ for parallel feature extraction across 200k+ dataset rows.
        \item Storage: SSD with at least 15 GB of free space for datasets and generated CSV files.
    \end{itemize}
\end{itemize}

\section{Non-Functional Requirements}

Non-functional requirements define the performance, scalability, security, and reliability characteristics of the SafeSurf system.

\subsection{SafeSurf System Constraints and Quality Attributes}

\subsubsection*{Performance}

\begin{itemize}
    \item The system is designed to provide responses within an average time of less than 500 milliseconds under normal operating conditions.
    
    \item It supports real-time processing of URLs with optimized feature extraction and prediction pipelines.
    
    \item External network operations such as DNS and WHOIS queries are controlled using timeout limits (e.g., 15 seconds) to prevent blocking.
\end{itemize}

\subsubsection*{Scalability}

\begin{itemize}
    \item The system uses concurrency control mechanisms such as semaphores to efficiently manage multiple simultaneous requests.
    
    \item It supports horizontal scaling through containerized deployment and cloud infrastructure.
    
    \item Future enhancements include caching mechanisms (e.g., Redis) to improve response time for frequently accessed URLs.
\end{itemize}

\subsubsection*{Availability}

\begin{itemize}
    \item The system targets an uptime of approximately 99.8\% and is designed to handle partial failures gracefully.
    
    \item Failures from external services such as DNS resolution, WHOIS lookups, or SSL verification are handled without crashing the system.
    
    \item The system supports dynamic model reloading (\texttt{/reload-model}) without requiring a full restart, ensuring continuous availability.
\end{itemize}

\subsubsection*{Security}

\begin{itemize}
    \item Input validation and sanitization are applied to all incoming URLs to prevent malformed or malicious inputs.
    
    \item Request size limits are enforced using middleware (e.g., \texttt{MAX\_REQUEST\_BODY\_BYTES}) to prevent large payload attacks.
    
    \item All communication between client and server is secured using HTTPS.
    
    \item Cross-Origin Resource Sharing (CORS) policies are configured to restrict access to trusted sources only.
\end{itemize}

\subsubsection*{Reliability}

\begin{itemize}
    \item The system maintains a prediction accuracy greater than 97\% based on trained models.
    
    \item It ensures consistent outputs across different datasets and input variations.
    
    \item The system is designed to remain stable under varying workloads and concurrent usage conditions.
\end{itemize}

\subsubsection*{Usability}

\begin{itemize}
    \item The interface is designed to be easy to understand and accessible to non-technical users.
    
    \item Results are presented in a clear and concise format without exposing internal technical details.
    
    \item The Chrome extension operates automatically, requiring minimal user interaction.
\end{itemize}

\subsubsection*{Maintainability}

\begin{itemize}
    \item The system follows a clean and modular code structure separating API routes, feature extraction logic, and model inference.
    
    \item Configuration settings are managed using environment variables (\texttt{.env}) for flexibility.
    
    \item Logging and monitoring mechanisms can be integrated to track system performance and errors.
\end{itemize}

\subsubsection*{Deployment \& Supportability}

\begin{itemize}
    \item The application is containerized using Docker for consistent deployment across environments.
    
    \item It is compatible with cloud platforms and supports scaling based on demand.
    
    \item The system is designed to integrate with CI/CD pipelines for automated deployment and updates.
\end{itemize}